\chapter{Discussion} \label{chap:discussion}
% Discussion helps to make your paper memorable - consolidate (summarize) your data to make it easier to remember. Briefly summarize your Methods (what you did) and Results (what you got and what you learnt).

% Skeleton:
% \begin{itemize}
%     \item Recap your Methods
%     \item Recap your Results
%     \item List other researchers’ reports using similar methods
%     \item Compare you results to theirs
%     \item \begin{itemize}
%         \item List similarities and differences and/or
%         \item Make a prediction and/or
%         \item Describe a relevant empirical rule
%     \end{itemize}
% \end{itemize}

The results demonstrate that successful cross-sensor transfer from Sentinel-2 to VEN$\mu$S is possible across all investigated settings. However, the performance strongly depends on both the transfer learning strategy and the normalization approach applied. The following sections discuss the main factors influencing these results, compare the findings with related work, and identify limitations and directions for future research. 

\section{Effect of Normalization Strategy}
\label{sec:discussion_normalization}

Normalization strategy turned out to be one of the most influential variables in this study, with the influence varying across transfer learning settings. In the transductive setting, the choice of normalization strategy was the critical factor between total model failure and strong performance. Both raw and fixed Z-score with Sentinel-2 statistics lead to nearly zero cloud detection score, while dynamic Z-score achieved an mF1 of 0.919. This shows that while spectral alignment is necessary for cross-sensor transfer, it is not sufficient if absolute value ranges of the two datasets differ.

One reason for this difference lies in the difference in product level between the two datasets. The Sentinel-2 training data consists of Level~1C top-of-atmosphere reflectances, while VEN$\mu$S Level~2 products provide atmospherically corrected surface reflectances, resulting in systematically different value ranges across sensors. Raw strategy fails because the model receives VEN$\mu$S reflectance values that are systematically lower than those seen during Sentinel-2 training, placing inputs outside the value range the model was trained for. Fixed Z-score with Sentinel-2 statistics fails for the same reason. While it normalizes the data, it does so using statistics derived from the Sentinel-2 training distribution, meaning VEN$\mu$S values still fall far outside the normalized range seen during training. Fixed Z-score with VEN$\mu$S statistics partially addresses this by mapping the target data to its own distribution, but Z-score normalization only aligns the mean and standard deviation and does not account for differences in histogram shape arising from the distinct sensor characteristics and product levels of the two datasets. Dynamic Z-score effectively bypasses this problem by normalizing each tile independently to zero mean and unit standard deviation, ensuring that the model receives inputs in a consistent range regardless of the sensor or product level.

In the inductive setting, the normalization pattern shifts. Fixed Z-score with VEN$\mu$S statistics consistently outperforms or matches dynamic Z-score across all fine-tuning strategies, though the gap between the two is small in most cases and becomes negligible under full fine-tuning. This suggests that when sufficient target domain adaptation is applied through fine-tuning, the model is able to compensate for residual radiometric differences, making the choice between these two strategies less critical. The from-scratch results are somewhat different, with raw normalization performing best. Since all VEN$\mu$S scenes cover a small, spectrally consistent region of Israel, the radiometric properties are similar across tiles. In this setting, per-tile normalization likely removes useful spectral information rather than harmful radiometric variation, which explains why dynamic Z-score underperforms raw normalization in the absence of cross-sensor domain shift.

Taken together, these findings suggest that the optimal normalization strategy depends on the transfer learning setting and the geographic diversity of the dataset. If a single strategy is to be selected for general use, dynamic Z-score normalization is the most robust choice, as it performs well across all experiments.

\section{Transfer Learning Strategies}
\label{sec:discussion_tl}

The strong performance achieved by output layer fine-tuning suggests that the feature representations learned by the encoder and decoder on Sentinel-2 data are largely transferable to the VEN$\mu$S domain. This is consistent with the spectral similarity between the two sensors in the shared bands. The pretrained network appears to have learned cloud-related spectral and spatial patterns that remain meaningful across these sensors, requiring only a remapping of the final class boundaries to adapt to the target domain.

Frozen encoder fine-tuning outperforming full fine-tuning is consistent with the findings of Yosinski et al.~\cite{yosinski2014}. Yosinski et al. suggest that lower layers learn more general and transferable features while higher layers become increasingly specialized to the source domain. Fine-tuning the decoder while keeping the encoder frozen therefore represents the optimal trade-off, where the decoder adapts to the spatial and radiometric properties of the target domain while the encoder representations are preserved. The small performance difference between frozen encoder and full fine-tuning further suggests that updating the encoder weights during fine-tuning does not provide additional benefit for this sensor pair.

An interesting observation relates to the potential role of shadow class knowledge in the pretrained model. The Sentinel-2 pretrained model was trained with explicit cloud shadow annotations, requiring the network to learn to distinguish shadows from other dark surfaces such as water bodies and shaded terrain. Even though the output layer is replaced and evaluation is performed on two classes only, this shadow-discriminative knowledge may be retained in the encoder representations and could provide an implicit advantage for cloud detection in the target domain. This hypothesis could be tested in future work by repeating all experiments with the cloud shadow class merged into the clear class in the Sentinel-2 training data, producing a 2-class source model that lacks explicit shadow knowledge, and comparing its performance against the results of the present study. If the transfer learning models trained without shadow knowledge underperform those trained with it, this would confirm that the shadow class knowledge encoded in the pretrained model provides an advantage for cloud detection in the target domain.

An important finding of the transductive setting is that the 3-class output of the Sentinel-2 pretrained model, including cloud shadow predictions, was successfully transferred to VEN$\mu$S imagery without any modification to the model weights. Visual inspection of prediction results using dynamic Z-score normalization showed that shadows were detected in plausible locations.

Visual inspection of the prediction examples provides additional insight into model behaviour across different land cover types. In the transductive setting, dark surfaces such as water bodies are sometimes misclassified as shadows, and bright desert surfaces are occasionally misclassified as clouds. In the inductive setting, predictions are generally accurate across all land cover types. A minor issue is observed above urban areas, where bright white rooftops are occasionally misclassified as cloud. Visual comparison also suggests that in some locations the model predictions may even be more accurate than the ground truth annotations.

A lower learning rate of $10^{-4}$ was evaluated for selected fine-tuning experiments as an alternative to $10^{-3}$ used by default. Despite being commonly recommended for fine-tuning to avoid catastrophic forgetting \cite{goodfellow2016}, the lower learning rate did not improve performance and resulted in significantly longer training times.

\section{Comparison with Related Work}
\label{sec:discussion_related}

As introduced in Section~\ref{sec:cnn_cloud}, Mateo-García et al.~\cite{garcia2020} investigated cross-sensor transfer learning between Landsat 8 and Proba-V, providing the closest comparable study to this work. Their finding that models can be trained before the target satellite is launched aligns with the assumption adopted here. In the inductive setting, they found no significant improvement with fine-tuning, which contrasts with the findings of this study. Also both sensors in their study used top-of-atmosphere reflectances as input, meaning they did not encounter the product level mismatch present in this study, which may explain why normalization strategy was not investigated in their work. Methodologically, they used spectral response functions to align and weight band contributions, while this study uses simple central wavelength matching.

Wright et al.~\cite{wright2025}, mentioned in Section~\ref{sec:cnn_cloud}, demonstrated that dynamic Z-score normalization enables cross-sensor generalization without any target domain adaptation. This study confirms that dynamic Z-score normalization is effective for cross-sensor transfer, identifying dynamic Z-score as the best strategy in the transductive setting and the most robust choice overall. However, while Wright et al.\ focused on building a single robust model for multiple sensors, this study investigates fine-tuning strategies for targeted adaptation to a specific sensor.


\section{Limitations and Future Work}
\label{sec:discussion_limitations}

Several limitations of this study should be acknowledged. Both the source and target models in this study are regional. The Sentinel-2 training data covers the Middle East and North Africa region, and the VEN$\mu$S dataset covers only part of Israel. While the geographic alignment between the two datasets likely benefits transfer learning performance, it also means that the results may not generalize to other geographic regions. Nevertheless, it should be noted that the presence of desert surfaces represents a particularly challenging case for cloud detection. It would be worth investigating how the same transfer learning framework performs in a more global setting. For example using a globally trained Sentinel-2 model trained on the full CloudSEN12+ dataset as the source. However, this would be computationally expensive given the size of the full dataset. The regional scope was a deliberate choice in this study, as a globally trained source model would likely increase the domain shift in the transductive setting. Whether this increased domain shift could still be compensated by dynamic Z-score normalization alone remains an open question.

The annotation quality of the CloudSEN12+ dataset represents an additional source of uncertainty. The dataset is described as expert-labeled, yet a manual quality review was necessary to remove patches with clearly incorrect or incomplete annotations. Therefore, the annotation quality is not as high as the authors state. Visual inspection of the prediction examples suggests that in some locations the model predictions appear more accurate than the ground truth annotations, indicating that reported metrics of the Sentinel-2 pretrained model may underestimate its true performance.

This study focuses on a single source-target sensor pair, and the generalizability of the findings to other cross-sensor transfer scenarios involving different sensors or spatial resolutions remains to be investigated. A related experiment worth investigating is the performance of the same framework using only bands available on most optical sensors, such as RGB and near-infrared, which would provide insight into the minimum spectral configuration required for effective cross-sensor transfer.

Several further methodological extensions could be explored. More advanced domain adaptation techniques, including adversarial approaches, could be investigated as alternatives to input-level normalization. Joint training, which combines labeled data from both source and target domains, represents another inductive transfer learning strategy not investigated in this study. Future work could also investigate spectral adaptation approaches that remove the requirement for spectrally equivalent bands. Finally, more recent transformer-based architectures could be explored as alternatives to the U-Net used in this study, potentially offering improved performance.
